\documentclass[12pt,letterpaper]{article}

\usepackage[T1]{fontenc}
\usepackage[utf8]{inputenc}
\usepackage{lmodern}
\usepackage[margin=1in]{geometry}
\usepackage{microtype}
\usepackage{hyperref}
\usepackage{setspace}

\hypersetup{
  colorlinks=true,
  linkcolor=black,
  urlcolor=black,
  citecolor=black
}

\title{Derived Agent Personalities: Grounding, Ownership, and Homogeneity in LLM-Based Builders}
\author{Anil Kumar\\Madras School of Economics}
\date{\today}

\begin{document}

\maketitle

\begin{abstract}
\noindent
Large language model agents are increasingly used as general-purpose coding and creative partners, yet default system configurations produce homogeneous outputs: similar tone, similar assumptions, and interchangeable ``helpful assistant'' behavior.
We characterize this failure mode as \emph{agent output homogeneity}: technically competent responses that lack stable, differentiated identity.
Hand-written persona prompts partially address this problem but do not scale, are difficult for non-expert users to author, and rarely connect identity to persistent ownership or provenance.

This thesis investigates \textbf{derived agent personalities}: structured identity artifacts generated from external, verifiable data and deployed as portable system prompts across agent interfaces.
We formalize personality as a multi-component artifact comprising voice, role, behavioral constraints, strengths, blind spots, and a canonical system prompt, rather than as ad hoc chat styling.
Our approach separates \textbf{identity generation} (deterministic feature extraction plus constrained language-model synthesis) from \textbf{identity deployment} (ownership-gated access and export into third-party agent environments).

We implement and evaluate this framework through Normie Works, a production system that derives coworker personas from on-chain NFT attributes, including visual, trait-level, and edit-history signals, and gates full personality access via cryptographic wallet authentication.
Public surfaces expose a shareable character card; private surfaces deliver the complete system prompt and an in-app conversational agent.
We study (1)~\textbf{distinctiveness} across tokens, (2)~\textbf{consistency} of voice within and across sessions, (3)~\textbf{grounding fidelity} to source attributes, and (4)~\textbf{downstream utility} when holders import personas into vibecoding workflows versus generic baselines.

Our results suggest that personality is most effectively treated not as model fine-tuning alone, but as an \textbf{identity layer} in the agent stack: generated, provenance-linked, and portable.
This work contributes a pipeline for grounded persona synthesis, an ownership model for holder-specific agent identity, and empirical evidence on whether derived personalities reduce homogeneity and improve user-agent collaboration in semi-technical builder settings.
We discuss implications for NFT holder utility, creator tooling, and the design of future agent platforms where character is a first-class primitive.
\end{abstract}

\section*{Research Question}
Can agent personality be treated as a grounded, ownership-linked artifact, generated from verifiable data and exported across tools, rather than as generic system-prompt prose?

\section*{Planned Contributions}
\begin{enumerate}
  \item A formal decomposition of agent personality into generatable, evaluable components.
  \item A pipeline combining deterministic feature extraction and temperature-0 LLM synthesis with version-keyed caching.
  \item An ownership-gated release model linking persona access to on-chain holder identity.
  \item Empirical analysis of distinctiveness, consistency, grounding fidelity, and builder utility versus generic baselines.
\end{enumerate}

\section*{System Reference}
Production implementation: \url{https://normies.sandpark.co}

\end{document}
